\subsection{Closed-sector contractions and coherent rephasing}

\begin{definition}[Closed sector tensors]%
    \label{def:mpdo_closed_sector_tensors}
    \lean{MPOTensor.closedSectorL}
    \lean{MPOTensor.closedSectorR}
    \leanok
    \uses{def:mpdo_sector_tensors}
    Tracing the two physical legs of a sector tensor leaves a vector over
    the virtual bond:
    \begin{align}
        |l_k)_{\gamma}&:=\tr\left[(l_k)_{\gamma}\right], \notag\\
        (r_k|_{\gamma}&:=\tr\left[(r_k)_{\gamma}\right]. \notag
    \end{align}
    These are the closed tensors of
    \cite[Appendix~C.2, lines~1473--1477]{Cirac2016MPDO_arXiv}.
\end{definition}

\begin{definition}[Closed-sector pairing operator]
    \label{def:mpdo_closed_sector_pairing_operator}
    \lean{MPOTensor.closedSectorPairingOperator}
    \leanok
    \uses{def:mpdo_closed_sector_tensors}
    For sector tensors $(l_k)_\beta$ and $(r_k)_\alpha$, close their physical
    legs and set
    \begin{align}
        S&:=\sum_k |l_k)(r_k|. \label{eq:rfp_pairing_operator}
    \end{align}
    This is the pairing operator displayed in
    \cite[Appendix~C.2, lines~1473--1493]{Cirac2016MPDO_arXiv}.
\end{definition}

\begin{theorem}[Sector trace pairing]
    \label{thm:mpdo_sector_trace_pairing}
    \lean{MPOTensor.trace_sectorEta}
    \leanok
    \uses{def:mpdo_sector_eta, def:mpdo_closed_sector_tensors}
    For all sectors $k,h$, the trace of the neighboring operator is the
    pairing of the closed sector tensors:
    $T_{k,h}=\tr(\eta_{k,h})=(r_k|l_h)$.
    This is \cite[Appendix~C.2, lines~1452--1455
    and~1478--1481]{Cirac2016MPDO_arXiv}.
\end{theorem}

\begin{proof}
    \leanok
    Expanding the definition of the neighboring operator and applying
    $\tr(A\otimes B)=\tr(A) \tr(B)$ gives
    \begin{align}
        \tr(\eta_{k,h})
        &=\sum_{\gamma}\tr((r_k)_{\gamma}\otimes(l_h)_{\gamma})
          =\sum_{\gamma}\tr[(r_k)_{\gamma}]\,\tr[(l_h)_{\gamma}]
          =(r_k|l_h). \notag
    \end{align}
\end{proof}

\begin{definition}[Concrete closed-sector trace matrix]
    \label{def:mpdo_concrete_closed_sector_trace_matrix}
    \lean{MPOTensor.closedSectorTraceMatrix}
    \lean{MPOTensor.closedSectorTraceMatrix_apply}
    \leanok
    \uses{def:mpdo_closed_sector_tensors, thm:mpdo_sector_trace_pairing}
    For the inverse-map sector tensors, define the complex matrix
    $T_{k,h}:=(r_k|l_h)$. Equivalently,
    $T_{k,h}=\tr(\eta_{k,h})$.
    This is the trace matrix in
    \cite[Appendix~C.2, lines~1478--1481]{Cirac2016MPDO_arXiv}.
\end{definition}

\begin{theorem}[Closed-sector form of the physical-trace transfer]
    \label{thm:mpdo_closed_sector_physical_trace_transfer}
    \lean{MPOTensor.physTraceTransfer_eq_sum_closedSector}
    \leanok
    \uses{def:mpo_physical_trace_transfer,
        def:mpdo_closed_sector_pairing_operator}
    More generally, any sector factorization
    $U_B\kappa_{\beta,\alpha}U_B^\dagger
    =\bigoplus_k(l_k)_\beta\otimes(r_k)_\alpha$ gives the closed-sector
    pairing operator
    \begin{align}
        \mathcal T_{\cal K}&=\sum_k |l_k)(r_k|. \label{eq:rfp_trace_transfer}
    \end{align}
    This identifies the transfer matrix of
    Definition~\ref{def:mpo_physical_trace_transfer} with the operator on the
    left-hand side of the zero-correlation-length identity in
    \cite[Appendix~C.2, lines~1489--1493]{Cirac2016MPDO_arXiv}.
\end{theorem}

\begin{proof}
    \leanok
    Trace the sector factorization over the physical index. Unitary
    invariance gives
    $\tr(U_B\kappa_{\beta,\alpha}U_B^\dagger)
    =\tr(\kappa_{\beta,\alpha})$.
    The trace of the block diagonal is the sum of the block traces. In each
    sector,
    \begin{align}
        \tr((l_k)_\beta\otimes(r_k)_\alpha)
        &=\tr[(l_k)_\beta]\,\tr[(r_k)_\alpha]
          =[|l_k)(r_k|]_{\beta,\alpha}. \notag
    \end{align}
    Summing over $k$ gives (\ref{eq:rfp_trace_transfer}).
\end{proof}

\begin{corollary}[Physical-trace transfer from the inverse-map sectors]
    \label{cor:mpdo_concrete_closed_sector_physical_trace_transfer}
    \lean{MPOTensor.concrete_physTraceTransfer_eq_sum_closedSector}
    \leanok
    \uses{def:mpdo_injective_inverse_layer, def:mpdo_eta_structure,
        thm:mpdo_closed_sector_physical_trace_transfer}
    Let ${\cal K}$ be injective. Let $R$ and $\rho$ give a three-site
    closure of ${\cal K}$, choose a Hayashi decomposition of $\rho$, and fix
    virtual indices $\alpha_1,\beta_3$ such that
    $R_{\beta_3,\alpha_1}\ne0$. For the corresponding closed sector tensors,
    $\mathcal T_{\cal K}=\sum_k |l_k)(r_k|$.
\end{corollary}

\begin{proof}
    \leanok
    \uses{thm:mpdo_sector_factorization,
        thm:mpdo_closed_sector_physical_trace_transfer}
    The inverse-map sector tensors satisfy the sector factorization of
    Theorem~\ref{thm:mpdo_sector_factorization}. Apply
    Theorem~\ref{thm:mpdo_closed_sector_physical_trace_transfer}.
\end{proof}

\begin{theorem}[Normalized closed-sector pairing operator]
    \label{thm:mpdo_closed_sector_pairing_normalized_idempotent}
    \lean{MPOTensor.closedSector_operator_quasi_idempotent}
    \lean{MPOTensor.closedSector_operator_normalized_idempotent}
    \lean{MPOTensor.closedSector_operator_idempotent_of_physTraceTransfer_sq}
    \leanok
    \uses{def:mpo_source_zcl,
        def:mpdo_injective_inverse_layer, def:mpdo_eta_structure,
        cor:mpdo_concrete_closed_sector_physical_trace_transfer,
        lem:mpo_source_zcl_normalized_idempotent}
    Let ${\cal K}$ be injective. Let $R$ and $\rho$ give a three-site
    closure of ${\cal K}$, choose a Hayashi decomposition of $\rho$, and fix
    virtual indices $\alpha_1,\beta_3$ such that
    $R_{\beta_3,\alpha_1}\ne0$. Form the corresponding closed sector
    tensors $|l_k)$ and $(r_k|$. If ${\cal K}$ satisfies the positive-real
    up-to-scalar physical-trace relation, set
    $S:=\sum_k |l_k)(r_k|$. Then there is a real number $\lambda>0$ such that
    \begin{align}
        S^2&=\lambda S, \label{eq:rfp_pairing_quasi_idempotent}\\
        (\lambda^{-1}S)^2&=\lambda^{-1}S. \notag
    \end{align}
    If ${\cal K}$ satisfies the paper's literal source zero-correlation-length
    identity $\mathcal T_{\cal K}^2=\mathcal T_{\cal K}$, then the raw pairing
    operator itself satisfies $S^2=S$. This is the display preceding
    \cite[Appendix~C.2, Lemma~C.5]{Cirac2016MPDO_arXiv}. The first identity
    records the broader relation used in this development, which is invariant
    under positive real rescaling; the paper's identity is its $\lambda=1$
    specialization.
\end{theorem}

\begin{proof}
    \leanok
    \uses{cor:mpdo_concrete_closed_sector_physical_trace_transfer,
        lem:mpo_source_zcl_normalized_idempotent}
    Substitute
    $S=\mathcal T_{\cal K}$ from
    Corollary~\ref{cor:mpdo_concrete_closed_sector_physical_trace_transfer}
    into the defining relation
    $\mathcal T_{\cal K}^2=\lambda\mathcal T_{\cal K}$, and then apply
    Lemma~\ref{lem:mpo_source_zcl_normalized_idempotent} to obtain
    (\ref{eq:rfp_pairing_quasi_idempotent}) and its normalized form. Under
    literal idempotence, the same substitution gives $S^2=S$ directly.
\end{proof}

\begin{theorem}[Normalized closed-sector trace relations]
    \label{thm:mpdo_closed_sector_trace_normalized_pow_relation}
    \lean{MPOTensor.closedSectorTraceMatrix_normalized_relations}
    \leanok
    \uses{def:mpdo_concrete_closed_sector_trace_matrix,
        thm:mpdo_closed_sector_pairing_normalized_idempotent}
    Under the hypotheses of
    Theorem~\ref{thm:mpdo_closed_sector_pairing_normalized_idempotent},
    suppose that ${\cal K}$ satisfies the positive-real up-to-scalar
    physical-trace relation. There is a real number $\lambda>0$ such that, for
    $\widehat T:=\lambda^{-1}T$, one has
    \begin{align}
        \widehat T^2&=\widehat T^3. \label{eq:rfp_trace_square_cube}
    \end{align}
    For every positive integer $N$, one also has
    \begin{align}
        \tr(\widehat T^N)&=\tr(\widehat T). \label{eq:rfp_trace_powers}
    \end{align}
    For the literal $\lambda=1$ source zero-correlation-length identity, this
    specializes to the trace-power display in
    \cite[Appendix~C.2, lines~1490--1497]{Cirac2016MPDO_arXiv}. No idempotence
    or rank-one conclusion for $\widehat T$ is asserted.
\end{theorem}

\begin{proof}
    \leanok
    \uses{thm:mpdo_closed_sector_pairing_normalized_idempotent}
    Let $L$ be the matrix whose $h$th column is $|l_h)$, and let $Q$ be the
    matrix whose $k$th row is $(r_k|$. Then $S=LQ$ and $T=QL$. Place the
    normalization scalar on $L$. Since $\lambda^{-1}LQ$ is idempotent,
    associativity gives
    \begin{align}
        (Q\lambda^{-1}L)^2
        &=Q(\lambda^{-1}LQ)\lambda^{-1}L
          =Q(\lambda^{-1}LQ)^2\lambda^{-1}L
          =(Q\lambda^{-1}L)^3. \notag
    \end{align}
    Cyclicity of trace and the idempotence of $\lambda^{-1}LQ$ give
    \begin{align}
        \tr(Q\lambda^{-1}L)
        &=\tr(\lambda^{-1}LQ)
          =\tr((\lambda^{-1}LQ)^2)
          =\tr((Q\lambda^{-1}L)^2). \notag
    \end{align}
    By (\ref{eq:rfp_trace_square_cube}),
    $\widehat T^N=\widehat T^2$ for every $N\geq2$, which proves
    (\ref{eq:rfp_trace_powers}).
\end{proof}

\begin{definition}[Explicit neighboring $\eta$-operator data]
    \label{def:mpdo_explicit_eta_operators}
    \lean{MPOTensor.etaOperators}
    \lean{MPOTensor.ExplicitEtaOperators}
    \leanok
    \uses{def:mpdo_eta_structure}
    Fix a Hayashi decomposition witness $h_\eta$. An explicit neighboring
    $\eta$-family is a dependent family $(\eta_{k,h})$ indexed by sector pairs,
    where $\eta_{k,h}$ acts on the neighboring bond space
    $B_k^{R} \otimes B_h^{L}$. Requiring positivity for each pair gives the
    corresponding positive $\eta$-data.
\end{definition}

\begin{definition}[Neighboring operators from sector tensors]
    \label{def:mpdo_eta_of_sector_tensors}
    \lean{MPOTensor.etaOfSectorTensors}
    \leanok
    \uses{def:mpdo_sector_tensors, def:mpdo_explicit_eta_operators}
    Let $(l_k)_a$ and $(r_k)_a$ be sector tensors indexed by the horizontal
    virtual index $a$. For neighboring sectors $k,h$, define an operator on
    $B_k^R\otimes B_h^L$ by
    \begin{align}
        \eta_{k,h}((x_R,x_L),(y_R,y_L))
        &=\sum_a (r_k)_a(x_R,y_R)(l_h)_a(x_L,y_L).
        \label{eq:rfp_eta_entries}
    \end{align}
    Thus $\eta_{k,h}=r_k l_h$ is the contraction over the shared horizontal
    index in
    \cite[Appendix~C.2, equation~\textup{(etarl)}]{Cirac2016MPDO_arXiv}.
    Positivity is not part of this definition.
\end{definition}

\begin{lemma}[Entries of the neighboring operators]
    \label{lem:mpdo_eta_of_sector_tensors_apply}
    \lean{MPOTensor.etaOfSectorTensors_apply}
    \leanok
    \uses{def:mpdo_eta_of_sector_tensors}
    For every pair of sectors and every pair of matrix indices, the entry of
    $\eta_{k,h}$ is the contraction displayed
    in (\ref{eq:rfp_eta_entries}).
\end{lemma}

\begin{proof}
    \leanok
    This is immediate from the definition of $\eta_{k,h}$.
\end{proof}

Neighboring right and left sector tensors meet in the bond operator
$\eta_{k,h}=r_k l_h$:
\begin{center}
    % eta_{k,h} = r_k l_h: the neighboring bond operator is the
    % contraction of the right sector tensor r_k with the left sector
    % tensor l_h over their shared horizontal index a,
    % eta_{k,h}((x_R,x_L),(y_R,y_L))
    %   = sum_a (r_k)_a(x_R,y_R) (l_h)_a(x_L,y_L).
    % Ink: the bond r_k -- l_h is the summed index a; the up legs are
    % the row indices x_R, x_L and the down legs the column indices
    % y_R, y_L; the row ends are sealed because r_k and l_h carry no
    % further horizontal index.  Both sides carry the same boundary
    % signature
    % (virtual-west=0 | virtual-east=0 | physical-up=2 | physical-down=2).
    \tenkzkernel
    \begin{tenkz}[rows={wire}, cols=2, west=none, east=none]
        \tn[skin=pill, wide=2,
            ports={90@1:physical:$x_R$, 90@2:physical:$x_L$,
                   270@1:physical:$y_R$, 270@2:physical:$y_L$}]{\eta_{k,h}}
    \end{tenkz}
    $ = $
    \begin{tenkz}[rows={wire}, cols=2, west=none, east=none]
        \tn[label pos=w, ports={90:physical:$x_R$, 270:physical:$y_R$}]{r_k} &
        \tn[label pos=e, ports={90:physical:$x_L$, 270:physical:$y_L$}]{l_h}
        \tnmark[form=label, label pos=n]{midway (1,1) and (1,2)}{$a$}
    \end{tenkz}
\end{center}
This is the contraction in
\cite[Appendix~C.2, lines~1441--1445]{Cirac2016MPDO_arXiv}.

\begin{theorem}[Cyclic contraction of the sector tensors]
    \label{thm:mpdo_cyclic_sector_contraction}
    \lean{MPOTensor.sum_cyclic_sectorTensor_eq_prod_etaOfSectorTensors}
    \leanok
    \uses{def:mpdo_eta_of_sector_tensors}
    Let $N\geq1$ and let $k_0,\ldots,k_{N-1}$ be a cyclic sequence of
    sectors. Fix matrix indices in the left and right factor of every sector,
    and abbreviate
    \begin{align}
        L_n(a)&=(l_{k_n})_a(x_n^L,y_n^L), \notag\\
        R_n(a)&=(r_{k_n})_a(x_n^R,y_n^R). \notag
    \end{align}
    With all site labels read modulo $N$,
    \begin{align}
        \sum_{g_0,\ldots,g_{N-1}}
          \prod_{n=0}^{N-1}L_n(g_n)R_n(g_{n+1})
        &=\prod_{n=0}^{N-1}
          \eta_{k_n,k_{n+1}}
          ((x_n^R,x_{n+1}^L),(y_n^R,y_{n+1}^L)).
        \label{eq:rfp_cyclic_contraction}
    \end{align}
    This is the entrywise cyclic contraction at
    \cite[Appendix~C.2, lines~1446--1450]{Cirac2016MPDO_arXiv}.
\end{theorem}

\begin{proof}
    \leanok
    \uses{lem:mpdo_eta_of_sector_tensors_apply}
    Expand the product of the sums defining the neighboring operators and set
    $g'_n=g_{n+1}$. Cyclic invariance of the finite product and commutativity of
    scalar multiplication give
    \begin{align}
        \sum_{g_0,\ldots,g_{N-1}}
          \prod_{n=0}^{N-1}L_n(g_n)R_n(g_{n+1})
        &=\sum_{g'_0,\ldots,g'_{N-1}}
          \prod_{n=0}^{N-1}R_n(g'_n)L_{n+1}(g'_n) \notag\\
        &=\prod_{n=0}^{N-1}\sum_a R_n(a)L_{n+1}(a). \notag
    \end{align}
    The last expression is the product of the corresponding entries of
    $\eta_{k_n,k_{n+1}}$, as in (\ref{eq:rfp_cyclic_contraction}).
\end{proof}

\begin{theorem}[Closed sector word as neighboring contractions]
    \label{thm:mpdo_sector_word_eta_contraction}
    \lean{MPOTensor.trace_sector_word_eq_prod_etaOfSectorTensors}
    \leanok
    \uses{def:mpdo_eta_of_sector_tensors}
    Under the hypotheses of the preceding theorem, define
    $M_n(a,b)=L_n(a)R_n(b)$. Then
    \begin{align}
        \tr(M_0M_1\cdots M_{N-1})
        &=\prod_{n=0}^{N-1}
          \eta_{k_n,k_{n+1}}
          ((x_n^R,x_{n+1}^L),(y_n^R,y_{n+1}^L)).
        \label{eq:rfp_sector_word}
    \end{align}
\end{theorem}

\begin{figure}[ht]
    \centering
    \begin{align}
        \tr(M_{k_0}\cdots M_{k_{N-1}})
        &=\tr(\eta_{k_0,k_1}\otimes\cdots\otimes\eta_{k_{N-1},k_0}).
        \notag
    \end{align}
    \caption{Closing the horizontal indices of the factorized sector tensors
    pairs each right factor with the left factor at the next site. The cycle
    therefore becomes the product of the neighboring operators
    $\eta_{k_n,k_{n+1}}$; compare
    \cite[Appendix~C.2, lines~1446--1450]{Cirac2016MPDO_arXiv}.}
    \label{fig:mpdo_cyclic_eta_contraction}
\end{figure}

\begin{proof}
    \leanok
    \uses{thm:mpdo_cyclic_sector_contraction}
    Expanding the trace of the matrix word gives
    \begin{align}
        \tr(M_0M_1\cdots M_{N-1})
        &=\sum_{g_0,\ldots,g_{N-1}}
          \prod_{n=0}^{N-1}M_n(g_n,g_{n+1}). \notag
    \end{align}
    Substituting $M_n(a,b)=L_n(a)R_n(b)$ gives the cyclic sum
    in (\ref{eq:rfp_cyclic_contraction}); applying that identity yields
    (\ref{eq:rfp_sector_word}).
\end{proof}

\begin{definition}[Sector-adapted chain coordinates]
    \label{def:mpdo_sector_chain_coordinates}
    \lean{MPOTensor.SectorFiber}
    \lean{MPOTensor.SectorChainIndex}
    \lean{MPOTensor.sectorChainEquiv}
    \leanok
    \uses{def:mpdo_eta_structure}
    For a sector configuration $k=(k_0,\ldots,k_{N-1})$, let
    \begin{align}
        I_k&=\prod_{n=0}^{N-1}
          \bigl(\{0,\ldots,d_L(k_n)-1\}\times
                 \{0,\ldots,d_R(k_n)-1\}\bigr). \notag
    \end{align}
    Applying the one-site Hayashi decomposition at every site gives a
    bijection
    \begin{align}
        \Phi_N:\{0,\ldots,d-1\}^N
        &\longrightarrow
          \coprod_{k_0,\ldots,k_{N-1}} I_k. \notag
    \end{align}
    We write $\reindex_{\Phi_N}(M)$ for the matrix obtained by
    transporting both indices of $M$ along this bijection.
\end{definition}

\begin{definition}[Physical conjugation of an MPO tensor]
    \label{def:mpdo_conjugate_physical}
    \lean{MPOTensor.conjugatePhysical}
    \leanok
    \uses{def:mpo_tensor}
    Let $U$ be a matrix on the physical space. The physical conjugation of
    ${\cal K}$ by $U$ is the tensor $\widetilde{\cal K}$ defined by
    \begin{align}
        \widetilde\kappa_{\beta,\alpha}^{\,ij}
        &=(U\kappa_{\beta,\alpha}U^\dagger)_{ij}. \notag
    \end{align}
    Its horizontal bond indices are unchanged. When $U$ is unitary, this is
    the corresponding physical basis change.
\end{definition}

\begin{definition}[Cyclic tensor product of neighboring operators]
    \label{def:mpdo_cyclic_eta_tensor_product}
    \lean{MPOTensor.cyclicEtaTensorProduct}
    \leanok
    \uses{def:mpdo_sector_chain_coordinates,
        def:mpdo_explicit_eta_operators}
    Let $N\geq1$ and let $k$ be a sector configuration. Define the matrix
    $\Omega_k$ on $I_k$ by
    \begin{align}
        \Omega_k(x,y)
        &=\prod_{n=0}^{N-1}
          \eta_{k_n,k_{n+1}}
          ((x_n^R,x_{n+1}^L),(y_n^R,y_{n+1}^L)).
        \label{eq:rfp_cyclic_eta_product}
    \end{align}
    Thus
    $\Omega_k=\eta_{k_0,k_1}\otimes\cdots\otimes
    \eta_{k_{N-1},k_0}$ after the neighboring tensor factors are ordered
    cyclically.
\end{definition}

\begin{theorem}[A fixed sector block is the cyclic neighboring product]
    \label{thm:mpdo_fixed_sector_cyclic_eta}
    \lean{MPOTensor.mpo_submatrix_sector_eq_cyclicEtaTensorProduct}
    \leanok
    \uses{def:mpdo_conjugate_physical,
        def:mpdo_cyclic_eta_tensor_product,
        def:mpdo_eta_of_sector_tensors}
    Suppose that the physical slices in the Hayashi basis satisfy
    \begin{align}
        U_B\kappa_{\beta,\alpha}U_B^\dagger
        &=\bigoplus_q(l_q)_\beta\otimes(r_q)_\alpha. \notag
    \end{align}
    Let $\eta_{q,h}=r_q l_h$. For every $N\geq1$ and every sector
    configuration $k$, the corresponding diagonal block of
    $\reindex_{\Phi_N}(M_N(\widetilde{\cal K}))$ is $\Omega_k$.
\end{theorem}

\begin{proof}
    \leanok
    \uses{thm:mpdo_sector_word_eta_contraction}
    Fix $x,y\in I_k$ and expand the closed horizontal trace of the MPO word.
    The sector factorization writes each local entry as a product of entries
    of $l_{k_n}$ and $r_{k_n}$. Set
    \begin{align}
        M_n(\beta,\alpha)
        &=(l_{k_n})_\beta(x_n^L,y_n^L)
          (r_{k_n})_\alpha(x_n^R,y_n^R). \notag
    \end{align}
    Then
    \begin{align}
        \bigl[\reindex_{\Phi_N}
          (M_N(\widetilde{\cal K}))\bigr]_{(k,x),(k,y)}
        &=\tr(M_0M_1\cdots M_{N-1}) \notag\\
        &=\prod_{n=0}^{N-1}\eta_{k_n,k_{n+1}}
          ((x_n^R,x_{n+1}^L),(y_n^R,y_{n+1}^L))
          =\Omega_k(x,y). \notag
    \end{align}
    The second equality is
    Theorem~\ref{thm:mpdo_sector_word_eta_contraction}, and the last is
    (\ref{eq:rfp_cyclic_eta_product}).
\end{proof}

\begin{theorem}[Sector-adapted MPO decomposition]
    \label{thm:mpdo_sector_adapted_decomposition}
    \lean{MPOTensor.mpo_reindex_sectorChainEquiv_eq_blockDiagonal_cyclicEta}
    \leanok
    \uses{def:mpdo_sector_chain_coordinates,
        def:mpdo_conjugate_physical,
        def:mpdo_cyclic_eta_tensor_product,
        def:mpdo_eta_of_sector_tensors}
    Under the same sector factorization hypothesis, for every $N\geq1$,
    \begin{align}
        \reindex_{\Phi_N}
          (M_N(\widetilde{\cal K}))
        &=\bigoplus_{k_0,\ldots,k_{N-1}}
          \eta_{k_0,k_1}\otimes\eta_{k_1,k_2}\otimes\cdots\otimes
          \eta_{k_{N-1},k_0}.
        \label{eq:rfp_sector_decomposition}
    \end{align}
    This is the algebraic sector decomposition in
    \cite[Appendix~C.2, lines~1446--1450]{Cirac2016MPDO_arXiv}. It does not
    assert that the individual neighboring operators are positive.
\end{theorem}

\begin{proof}
    \leanok
    \uses{thm:mpdo_fixed_sector_cyclic_eta}
    If the row and column sector configurations coincide, the corresponding
    block is Theorem~\ref{thm:mpdo_fixed_sector_cyclic_eta}. If they differ,
    choose a site $n$ at which their labels $k_n$ and $h_n$ are unequal.
    Block diagonality gives
    $[U_B\kappa_{\beta,\alpha}U_B^\dagger]
    _{(k_n, \cdot),(h_n, \cdot)}=0$ for every pair of horizontal indices.
    This local factor occurs in every cyclic horizontal-index product, so
    every summand in the closed trace is zero, proving
    (\ref{eq:rfp_sector_decomposition}).
\end{proof}
